\documentclass[11pt]{book}
\usepackage{fontspec}
\setmainfont{Latin Modern Roman}
\setsansfont{Latin Modern Sans}
\usepackage{amsmath,amssymb,mathtools}
\usepackage{unicode-math}
\setmathfont{Latin Modern Math}
\usepackage{xurl}
\usepackage[hidelinks]{hyperref}
\newcommand{\sref}[2]{\hyperlink{source:#1:#2}{[S#2]}}
\newcommand{\cataloguescope}[1]{\paragraph{Mathematical question.}#1}
\begin{document}
\section{Selfridge-Wagstaff-Pomerance Prize Problem}\label{statement}

\subsection{Definitions and mathematical statement}

\paragraph{Shared notation.}$\mathbb N=\{1,2,\ldots\}$, and $\mathbb Z,\mathbb Q,\mathbb R,\mathbb C$ denote the integers, rationals, reals and complex numbers. Any other conventions are specified locally.


Let $(F_j)_{j\geq0}$ be the Fibonacci sequence with $F_0=0$, $F_1=1$ and $F_{j+2}=F_{j+1}+F_j$. A composite number is an integer greater than one that is not prime. Congruence $n\equiv3\text{ or }7\pmod{10}$ makes $n$ odd and coprime to $10$.
\paragraph{Statement recorded in the supplied source.}
Does there exist a composite integer $n$ satisfying
\[n\equiv3\text{ or }7\pmod{10},\qquad n\mid(2^n-2),\qquad n\mid F_{n+1}?\]
Find such an integer or prove that none exists. \sref{statement}{2}
\subsection{Short English statement}
Can a composite integer ending in 3 or 7 simultaneously pass the stated base-two divisibility test and divide the next Fibonacci number?
\subsection{Sources}
\begin{description}
\item[\hypertarget{source:statement:1}{S1}] ULAM.AI and the UnsolvedMath Contributors, \emph{UnsolvedMath: A Curated Collection of Open Mathematics Problems}, catalogue entry. CC BY 4.0. \url{https://huggingface.co/datasets/ulamai/UnsolvedMath}.
\item[\hypertarget{source:statement:2}{S2}] Andrew Shallue and Jonathan Webster, \emph{Fast tabulation of challenge pseudoprimes}, arXiv:1806.08697 (2018); Open Book Series 2 (2019), 411--428. Introduction and Section2, pages1--3. \url{https://arxiv.org/pdf/1806.08697}.
\end{description}
\label{end:statement}
\end{document}
